\begin{tikzpicture}[
  >={Latex[length=1.5mm,width=1.2mm]},
  every node/.style={font=\small},
  fline/.style={thin,ered,->}
]

% --- (a) t, p=p0 ---
\begin{scope}
  % Linee chiuse piccole (dipolo stretto)
  \foreach \s in {0.3,0.5}{
    \draw[fline,smooth] (0.15,0) .. controls (0.15,\s) and (-0.15,\s) .. (-0.15,0)
      .. controls (-0.15,-\s) and (0.15,-\s) .. (0.15,0);
    \draw[fline,smooth] (-0.15,0) .. controls (-0.15,\s) and (0.15,\s) .. (0.15,0)
      .. controls (0.15,-\s) and (-0.15,-\s) .. (-0.15,0);
  }
  \node[below] at (0,-0.8) {$t$, $\mathbf{p}=\mathbf{p}_0$};
  \node[above right] at (0.3,0.5) {(a)};
\end{scope}

% --- (b) t+T/4, p=0 ---
\begin{scope}[xshift=4cm]
  \fill (0,0) circle (1.5pt);
  \node[below] at (0,-0.8) {$t+\dfrac{T}{4}$, $\mathbf{p}=0$};
  \node[above right] at (0.3,0.5) {(b)};
\end{scope}

% --- (c) t+T/2, p=-p0 ---
\begin{scope}[yshift=-3cm]
  \foreach \s in {0.3,0.5,0.8,1.1}{
    \draw[fline,smooth,domain=-180:180,samples=40]
      plot ({0.3*\s*cos(\x)},{\s*sin(\x)});
  }
  \node[below] at (0,-1.3) {$t+\dfrac{T}{2}$, $\mathbf{p}=-\mathbf{p}_0$};
  \node[above right] at (0.5,1.0) {(c)};
\end{scope}

% --- (d) t+3T/4, p=0 ---
\begin{scope}[xshift=4cm,yshift=-3cm]
  \fill (0,0) circle (1.5pt);
  \foreach \s in {0.4,0.7,1.0}{
    \draw[fline,smooth,domain=-180:180,samples=40]
      plot ({0.4*\s*cos(\x)},{\s*sin(\x)});
  }
  \node[below] at (0,-1.3) {$t+\dfrac{3T}{4}$, $\mathbf{p}=0$};
  \node[above right] at (0.5,1.0) {(d)};
\end{scope}

% --- (e) t+T, p=p0 ---
\begin{scope}[yshift=-7cm]
  \foreach \s in {0.3,0.6,0.9,1.3,1.7}{
    \draw[fline,smooth,domain=-180:180,samples=50]
      plot ({0.3*\s*cos(\x)},{\s*sin(\x)});
  }
  \node[below] at (0,-2.0) {$t+T$\quad $\mathbf{p}=\mathbf{p}_0$};
  \node[above right] at (0.5,1.5) {(e)};
\end{scope}

% --- (f) Pattern completo ---
\begin{scope}[xshift=4cm,yshift=-7cm]
  \foreach \s in {0.3,0.6,0.9,1.3,1.7,2.2}{
    \draw[fline,smooth,domain=-180:180,samples=60]
      plot ({0.4*\s*cos(\x)},{\s*sin(\x)});
  }
  % Cerchi tratteggiati concentrici
  \foreach \r in {1.0,1.5,2.0}{
    \draw[dashed,thin,gray!60!black] (0,0) circle (\r);
  }
  \draw[<->,dashed,thin] (0,1.5) -- (0,2.0);
  \node[right] at (0.1,1.75) {$\lambda$};
  \node[above right] at (0.5,2.0) {(f)};
\end{scope}

\end{tikzpicture}
